\section{The Record Infrastructure}
\label{sec:infrastructure}

The front door can disclose only what the deployment has made observable. Ordinary discovery reaches existing matter within a party's possession, custody, or control.\lrfootnote{See \statcite{frcp34control}.} It cannot recreate an unrecorded version, an undocumented release decision, a complaint that reached no assigned owner, or stop authority that no one possessed. The deployment-record duty therefore begins as operating infrastructure and later becomes litigation infrastructure.

This Part defines the control sheet, interface, and separate record clocks; tests feasibility against FAA, FTC, FCRA, and AI Act comparators; and addresses minimization and defense use.

\subsection{The Deployment-Record Control Sheet}

Before release, the operator should adopt a written control sheet for each designated deployment, version family, operating context, and risk category. This is not a broad ethics charter. It links technical identifiers to organizational authority and locates the continuing records. Every component has an existing source inside a firm that already runs ordinary engineering governance: identity and scope come from the asset inventory and release management; risk and release owners from the approval workflow; the evaluation and decision record from whatever already documents a go or no-go; feedback and escalation from the incident queue; intervention from the runbook and the access-control matrix; provenance and retention from vendor contracts and the retention schedule. The sheet is a register of organizational decisions, not a capture of system behavior. The minimum record the sheet carries is a ceiling as much as a floor: it fixes what the record duty compels, and the Act commands nothing beyond it by default; Part IV.E develops the minimization discipline that ceiling implies. Table~\ref{tab:control-sheet} states the sheet's minimum components and the front-door function each performs.

\begin{longtable}{>{\raggedright\arraybackslash}p{0.25\textwidth} >{\raggedright\arraybackslash}p{0.33\textwidth} >{\raggedright\arraybackslash}p{0.32\textwidth}}
\caption{The Deployment-Record Control Sheet}\label{tab:control-sheet}\\
\toprule
Component & Required allocation & Front-door function \\
\midrule
\endfirsthead
\toprule
Component & Required allocation & Front-door function \\
\midrule
\endhead
Identity and scope & Application and model versions; purpose; affected population; data; retrieval; tools; credentials; vendors; limits; the designation's material-change threshold. & Defines the deployment and distinguishes versions, component changes, and independent modification. \\
Risk and release owners & Responsible executive and operational owner for each designated risk; approvers; authority and resources; conditions and review date. & Identifies who receives information and who can commission mitigation or withhold release. \\
Evaluation and decision record & Protocols, populations, thresholds, results, exceptions, rejected or conditional mitigations, and unresolved risk decisions. & Records the basis for authorization and supplies both knowledge and reasonableness evidence. \\
Feedback and escalation & Complaint and incident categories; monitoring; escalation channel; recipient; response deadline; decision and corrective action. & Preserves notice, separates report from response, and makes continuing control observable. \\
Intervention & Actors authorized to restrict, revoke credentials, update, pause, roll back, or terminate; triggers, response times, and tested procedures. & Establishes practical authority and whether a proposed intervention was available at the relevant time. \\
Provenance and retention & Version lineage; material configuration changes; event identifiers; record custodians; third-party locations; retention and hold procedures. & Enables identification, authentication, preservation, proportional production, and safe exit. \\
External-rights savings clause & Statement that internal allocation and contract labels do not waive affected-person rights or determine ultimate liability. & Prevents the record from operating as a private exculpatory agreement. \\
\bottomrule
\end{longtable}

The record should identify roles and offices as well as names. Personnel change; statutory responsibility should survive turnover. A deployment ID links the approved control sheet to the live configuration, and a versioned amendment records material change. The operator can then answer a front-door notice by retrieving a defined record rather than assembling an organizational history after the event.

The material-change threshold is especially important for learning and composite systems. A model update can alter capability while leaving the application version unchanged. A new retrieval source, memory setting, tool credential, ranking threshold, user population, or downstream purpose can change risk without changing model weights. The control sheet applies the material-change rule fixed by the designation: it records which changes required renewed evaluation or approval and which entered a lower-burden change log. A provider and operator can use different internal systems so long as the deployment record links them by version, time, and control lane.

A completed form is only the beginning of proof. The operative questions are whether signals reached people with authority and resources, interventions worked, and material changes triggered review. The form's legal effect is the one Part II.C and Model Act \S~8 assign: it creates no merits inference.\lrfootnote{See \bluecite[\S~16(a)]{restatementPhysicalEmotional2010}; \bluecite[\S~4(b)]{restatementProducts1998}. Jurisdictions may assign statutory compliance a different effect.}

\subsection{The Record Has a Standard Interface}

The minimum record should operate as a standard interface across the deployment rather than as a mandate to centralize every underlying artifact. The visible operator keeps a closed control sheet: deployment ID; covered purpose and population; application and model versions; material components and suppliers; record custodians; risk and release owners; relevant contract rights; and the event schema used for consequential actions. Each control holder keeps the deeper materials assigned to its lane. Stable identifiers join the control sheet to those materials when a notice or regulator demand arrives.

Three identifiers do most of the work. A \term{deployment identifier} (deployment ID) distinguishes the covered use from the same model used elsewhere. A \term{component-version identifier} ties model, application, retrieval, tool, policy, and configuration changes to a time interval. An \term{event identifier} connects an affected person or consequential decision to the versions and records that governed it; forwarded to covered control holders with the notice under Model Act \S~4, it serves as the ``common event identifier.'' Together, the three identifiers form a \term{correlation identifier}---a value propagated along the call path and bound to a fixed set of legal fields, a pattern engineering already runs as propagated trace context; the designation fixes which legal fields ride on it.\lrfootnote{See W3C, Trace Context (Level 1), W3C Recommendation (Nov. 23, 2021), \url{https://www.w3.org/TR/trace-context/} (standardizing the traceparent and tracestate headers that propagate a fixed trace-context field set across service boundaries, and the default propagation format of existing telemetry stacks).} The event identifier can be pseudonymous and separately linked to personal data. The system therefore retrieves the relevant slice without exposing unrelated users or requiring a later investigator to infer lineage from timestamps alone.

The interface should use a sectoral data dictionary. ``Model version'' can otherwise mean a marketing name, weight snapshot, API alias, application release, policy bundle, or customer configuration. The designation should define the fields needed to reproduce the legally relevant identity: supplier identifier; immutable or otherwise auditable version reference; activation window; configuration owner; material change; downstream integration; output type; and custodian. An enterprise may keep its own engineering format, but the first answer translates that format into the defined fields and identifies the source system from which each field came.

Integrity requires lineage rather than a self-authenticating compliance certificate. A control sheet should record who created and approved it, the date, later amendments, and the source records supporting the entry. Automated logs can carry hashes, signatures, or append-only controls where appropriate; lower-technology sectors can use ordinary version control and attestation. The legal function is the same: preserve a contestable trail from the live event to the record, not establish conclusively that every listed control operated as designed.

Two requirements are expressly excluded. The Act does not require write-once storage, and it does not require cryptographic notarization of entries. Where the designation and sectoral practice impose nothing more, ordinary version control plus an approval signature satisfies the integrity function. The duty is no wider than the record: it asks that a field, once recorded, remain attributable to a time, an author, and a source---not that the underlying system be instrumented.

Supplier contracts implement the interface. A visible operator should receive current component and version identifiers, material-change notice, an event-resolution method, a responsive agent, and a right to retrieve the minimum supplier fields on demand. The supplier should receive the operator's deployment identifier, designated purpose, configured population, material customer settings, and preservation contact. Acquisition, subcontracting, model substitution, and service termination should carry successor and export duties for the retained period. These terms keep the record reachable when corporate and technical boundaries change after deployment. Recent work on the AI procurement record reaches a closely related diagnosis: contracts can separate evidence, operational control, economic risk, and public accountability.\lrfootnote{See \bluecite{maxwellProcurementRecord2026} (abstract) (treating the procurement file as a governance record). The front door adds a nonwaivable affected-person-facing answer and event-specific enforcement path.}

The architecture controls cost by keeping routine control-sheet fields separate from retrieval after a defined event. Release and asset-management systems can supply version and supplier identity. Transaction, decision, and incident systems can supply event fields. Deeper evaluation, complaint, and source materials stay with their custodians unless a risk-specific request justifies access. Standardization moves expense toward one-time design and retrieval by identifier, away from repeated forensic reconstruction across every dispute.

Quality assurance can also be bounded. The operator periodically reconciles the live deployment against the control sheet, tests whether an event identifier resolves to the governing versions, confirms that supplier contacts and retrieval rights remain usable, and records exceptions. A regulator can audit samples without obtaining every event. The check asks whether the front door will open when needed: can the operator identify what ran, locate the holder, preserve the slice, and exercise the retrieval path?

\subsection{Identity, Feedback, and Event Records Have Different Clocks}

A usable infrastructure contains three linked records.

The \term{identity record} changes when the deployment changes. It contains the versioned architecture, purpose, population, component suppliers, material configuration, data and tool interfaces, owners, and intervention rights. Its retention period should cover the ordinary period in which a claim involving that deployment can arise, with sector-specific extensions for minors, latent physical injury, or continuing public entitlements.

The \term{feedback record} accumulates during operation. It contains defined complaints, evaluations, incident classifications, monitoring signals, escalations, decisions, and corrective actions. The designation should specify which events require individual retention, which may be aggregated, and which content should be minimized. A service can record that a crisis detector triggered, which policy and version governed, and what intervention followed without preserving every unrelated intimate conversation indefinitely.

The \term{event record} freezes a bounded slice when a consequential action or identified incident occurs. In employment it may contain the application, input fields, model and configuration versions, score or recommendation, threshold, downstream human action, notice, and review. In benefits it may contain the assessment data, calculation version, output, governing policy, explanation, appeal, and correction. In a companion deployment it may contain the account and age status, relevant interaction window, governing relational and crisis policies, detections, interventions, warnings, and account restrictions. The event record anchors causal sequence while leaving unrelated people and periods outside production.

These clocks improve privacy and accuracy simultaneously. A blanket command to retain all interactions invites excess. A deployment ID, event schema, and material-change log preserve the facts needed to identify the relevant window. Sectoral lawmakers can then choose periods and minimization rules that fit the interest rather than defaulting to permanent surveillance.

Retention should follow the legal and factual life of the interest. Employment and credit records can track the limitations and administrative-charge periods governing the designated decision, with a defined post-notice extension. A public-benefits event may need to persist through reconsideration, appeal, and the period for judicial review. A clinical deployment may require longer identity and validation lineage where injury can appear later, while still minimizing raw patient content under health-information rules. A relational service involving a minor can preserve crisis and intervention events for a longer protected period without retaining every ordinary conversation. The designation should state the clock, trigger, and lawful deletion rule rather than asking operators to guess from a generic command to keep ``AI data.''

Different clocks also prevent a change from overwriting history. A current model card cannot establish what governed an earlier event; a current contract cannot prove who controlled retrieval before an acquisition; a revised policy cannot show which warning appeared to an earlier user. The identity record freezes each material interval. The event record points to that interval. The feedback record then shows what the organization learned and changed between intervals. This structure turns evolution from an obstacle into an auditable sequence.

Responsibility for legal holds should also appear in the record. The verified notice described in Part II identifies the event, period, user or decision, and risk. The operator's hold reaches the matching identity, feedback, and event materials and gives identified control holders a common reference. If a record no longer exists, the response identifies the deletion rule, date, custodian, and possible substitute. That information permits a court to distinguish ordinary expiration, negligent noncompliance, and intentional concealment.

\subsection{Existing Regimes Demonstrate Administrability}

The proposal borrows implementation forms, not liability conclusions. Four regimes perform distinct comparator roles: organizational record creation, regulator-defined information demands, a private path from visible decisionmaker to upstream file holder, and a supranational access right without a litigation-facing record. The first three comparators are national regimes---an architecture offered as a workaround for the absence of central authority would not prove its own feasibility out of central authority's own instruments---while the fourth prices what a unified instrument adds and what it omits.

\subsubsection{FAA Safety Management Systems}

FAA Part 5 requires covered aviation organizations, on the applicable schedule, to develop and maintain an organization-wide safety management system. The rules assign an accountable executive, require management authority to accept safety risk, specify hazard identification and risk control, require assurance and corrective action, and preserve records. For specified organizations, the system includes confidential employee reporting without fear of reprisal.\lrfootnote{See \statcite{faaSMSPart5}.}

Part 5 demonstrates four institutional facts. A safety system can attach to an organization rather than a single component. Named executive and operational authority can coexist. Continuing assessment and correction can be recorded as operating requirements. Retention can be specified without waiting for a lawsuit. AI deployments involve different technologies and protected interests, so the aviation standard does not define reasonable AI care. It demonstrates an administrable precedent for making authority, feedback, intervention, and records a mandatory operating system rather than an aspirational policy.

Aviation reporting supplies a related design lesson. Candid incident and near-miss information has prevention value, while a system that treats every report as a confession encourages silence. The Aviation Safety Reporting Program uses third-party reporting and a qualified, nonpunitive enforcement policy rather than blanket immunity.\lrfootnote{See \bluecite[\P\P~2, 6.2, 8.1--8.3, 12.3]{aviationReporting}.} The narrower lesson here is to distinguish a report from the enterprise's response, while leaving any consequences for falsification, concealment, retaliation, or unauthorized bypass to governing law.

\subsubsection{The FTC Companion-Chatbot Orders}

The Federal Trade Commission's 2025 section 6(b) study shows that AI-deployment records can be specified with considerably more precision than ``produce the algorithm.'' The Commission ordered named firms to identify organizational roles, predeployment standards and thresholds, deployment approvers, considered and rejected mitigations, post-deployment monitoring, complaints, employee concerns, escalation, and response.\lrfootnote{See \bluecite[Definitions \& Instructions; specs. 3, 11--17]{ftcCompanionOrder2025}. The order gathers information; it is not a finding that a recipient violated law.}

Those specifications map directly onto the front-door functions. Organizational roles identify the control lane. Thresholds and approvers define release authority. Rejected mitigations record a decision without deciding whether it was unreasonable. Monitoring, complaints, employee concerns, escalation, and response create a time-stamped feedback path. A sectoral statute can standardize the same categories prospectively and permit a court to request only the subset connected to a verified risk.

NIST's AI Risk Management Framework can supply terminology for lifecycle roles, executive responsibility, third-party components, testing, feedback, monitoring, and decommissioning.\lrfootnote{See \bluecite[GOVERN 1.7, 2.1, 2.3; MAP 3.5; MEASURE 2.5; MANAGE 4.1]{nistAIRMF2023}.} Its voluntary framework remains a vocabulary source. FAA rules and a sector-specific statute supply legal obligation; the FTC orders demonstrate that the responsive information can be operationally defined.

\subsubsection{The FCRA's Private-Actor First-Answer Sequence}

The Fair Credit Reporting Act supplies the closest domestic analogue for a private-actor gateway. In employment, a user generally must give the consumer a copy of the report and a rights summary before taking an adverse action based on that report.\lrfootnote{See \statcite{fcraEmploymentPreAdverse}.} After a covered adverse action, the visible decision user must notify the consumer, identify the consumer reporting agency that supplied the report, state that the agency did not make the decision, and give notice of the right to obtain a free report and dispute its accuracy.\lrfootnote{See \statcite{fcraAdverseAction}.} The agency, in turn, discloses the bounded consumer file, its sources, and recent recipients; a dispute triggers reinvestigation and a defined route to the furnisher.\lrfootnote{See \statcite{fcraDisclosureChain}.}

That sequence separates merits from access and the visible actor from the upstream information holder. Notice of adverse action does not establish that the report was inaccurate or that either entity violated a merits duty. It routes the affected person to a bounded file and correction process. The analogy concerns duties assigned to private actors, not a unitary private remedy: section 1681m(h)(8) excludes private civil actions for failures under section 1681m and assigns that section to public enforcement, while other FCRA provisions carry their own enforcement allocations.\lrfootnote{See \statcite{fcraNoPrivateAdverseAction}; \casecite[822--23]{perryFCRA2006}.} FCRA also applies only within its consumer-report architecture, generally presupposes an existing file, and supplies neither an AI deployment map nor a material-control test. Its structural lesson remains powerful: federal law can require the visible decision user to identify an upstream record holder while expressly declining to treat that holder as the decisionmaker. The front door adds the ex ante event and control record for legislatively designated deployments.

\subsubsection{A Supranational Comparator: Explanation Without a Litigation-Facing Record}

The European Union's unified regime completes the comparator set. The AI Act mandates technical documentation, automatic logging, and record retention for high-risk systems, but those duties run toward regulators and the provider's own file. Article 86 separately gives an affected person a right to a clear and meaningful explanation of the role of the AI system in the decision-making procedure and of the main elements of the decision taken.\lrfootnote{See \statcite{euAIAct2024}, art. 86 (excluding Annex III, point 2). Article 86 applies from August 2, 2026, under Article 113's general timetable; Regulation (EU) 2026/1744 of the European Parliament and of the Council of 8 July 2026 (Digital Omnibus on AI), art. 1(40), OJ L, 2026/1744, 24.7.2026, deferred the Chapter III documentation and logging duties for Annex III high-risk systems to December 2, 2027, without touching Article 86.} That is an explanation right. It is not an event receipt, a version identification, a custodian map, or a control allocation; reaching the underlying record still depends on an intermediary's discretion and on procedures that presuppose an identified target. The sequencing makes the point concrete: the explanation right is already in force while the documentation and logging duties on which any meaningful explanation would draw do not yet apply---an access right that arrives before the record duty has nothing to reach. The litigation-facing counterpart---the proposed AI Liability Directive's mechanism for court-ordered disclosure of evidence concerning high-risk systems---was withdrawn before enactment.\lrfootnote{Proposal for a Directive on adapting non-contractual civil liability rules to artificial intelligence (AI Liability Directive), art. 3, COM (2022) 496 final (Sept. 28, 2022) (court-ordered disclosure and preservation of evidence concerning high-risk AI systems); Commission Work Programme 2025, Annex IV, COM (2025) 45 final (Feb. 11, 2025) (announcing the intended withdrawal); Withdrawal of Commission Proposals, OJ C/2025/5423, 6.10.2025.} A unified regime thus supplies a regulator-facing record and an explanation right, but neither an event receipt nor a custodian map.

The same instrument also answers the speed form of the unified-regime objection, in the enacting legislature's own words. The Digital Omnibus deferred the core high-risk obligations well past their original dates---standalone Annex III obligations from August 2, 2026 to December 2, 2027, and Annex I obligations to August 2, 2028---and recital 2 of the amending Regulation states the reason: ``the delayed preparation of standards, which provide technical solutions for providers of high-risk AI systems to ensure compliance with their obligations under that Regulation, and the delayed establishment of the governance and the conformity assessment frameworks at national level has resulted in a compliance burden that is heavier than expected.''\lrfootnote{Regulation (EU) 2026/1744, recital 2, OJ L, 2026/1744, 24.7.2026. A recital states the legislature's reasons and aids interpretation; it does not itself bind. The recital attributes the heavier burden to the delayed standards and the delayed national governance and conformity-assessment frameworks, not to the intrinsic cost of the duties, and adds that stakeholder consultations revealed ``the need for additional measures to facilitate and provide clarification on implementation and compliance.'' \emph{Id.}} Unified enactment concentrates the decision; it cannot conjure the machinery that makes the decision operable. What slipped was not legislative will but implementation infrastructure---standards that did not yet exist and conformity-assessment bodies not yet established. The shape of that failure is this Article's thesis one level up: the binding constraint was not the absence of a rule but the absence of what the rule presupposed.

Read precisely, the recital distinguishes two kinds of ex ante duty. What proved unaffordable on schedule was an \term{approval architecture}---obligations discharged through harmonised standards and national conformity assessment, inert until both exist. The front door runs on a \term{record architecture} and depends on neither: the designation fixes the first answer's fields, in closed form, before the event, so what compliance requires is knowable on the day the statute passes, and there is no notified body, no conformity assessment, no CE marking, and no standards body on whose timetable the duty waits. That is why this Part's comparators are record regimes rather than conformity regimes. A description can be specified by a legislature in one sitting; a measure has to wait for someone to build the instrument.

Together, the comparators answer feasibility at the right level. FAA Part 5 shows prospective organizational records; the FTC orders show that deployment roles and events can be specified; FCRA shows a private-actor routing sequence across a divided decision stack; and the AI Act shows that even a unified supranational regime supplies a regulator-facing record and an explanation right without an event receipt or custodian map. For firms already capable of version control, release management, access logging, vendor management, incident response, and legal retention, the proposal connects those systems for designated deployments. Sectoral designations can provide simplified forms and regulator assistance where appropriate. None of these regimes requires a universal explanation of every internal model state.

\subsection{Recordkeeping Risks Require Minimization by Design}

Recordkeeping creates a distinct governance risk because recording can reshape the organization being observed. In a 2025 survey of one hundred practitioners, respondents described employee resistance, privacy tension, and uses extending beyond the accountability purpose; forty-seven percent reported records used to monitor or improve staff performance.\lrfootnote{See \artcite[554--56, 559--64]{chappidiAccountabilityCapture2025}.} A mandate framed as ``retain everything'' could expand employee and user surveillance, suppress candid internal discussion, preserve sensitive data without purpose, and reward litigation defensibility over safety value.

The front door answers that risk in its data structure. The operator's control sheet holds the minimum record rather than raw content by default. Identity, feedback, and event records use different clocks. The event schema records that a threshold fired, which policy governed, and what response followed; it does not require every conversation, keystroke, or unrelated user's data. Employee-performance monitoring is not a minimum field. Raw deliberative conversation, protected health information, and intimate or third-party content remain outside unless a sectoral designation expressly identifies a necessity and lawful protection. Automated capture must be limited to defined events, stated purposes, authorized recipients, and a deletion rule.

The organizational record also separates report from decision. A good-faith risk report records a signal, not an admission by the enterprise; the response record identifies who reviewed it, what information was available, and what action followed. That separation supports candor without making a report's existence automatic breach. Periodic sampling tests whether the record resolves an event without converting continuous worker observation into the measure of compliance. These specifications make minimization part of observability rather than a promise added after collection.

\subsection{The Record Is Also a Defense Asset}

The same record can defeat a claim. Version and configuration history may show that a feature was inactive, a downstream actor removed a safeguard, a warning appeared, an intervention fired, or the alleged controller lacked stop authority. Evaluations and incident sequences may establish reasonable preparation, intervening modification, user circumvention, or residual risk. They also keep the inquiry historical: later research may identify a question, and a later mitigation a possible control, without proving earlier knowledge, feasibility, or admissibility.\lrfootnote{See Fed. R. Evid. 407.} The front door prevents either side from asking a court to infer the deployment from the output alone.
